% !TEX program = xelatex
\documentclass[12pt,a4paper]{article}

% ── XeLaTeX: Unicode + system fonts ──────────────────────────────────────────
\usepackage{fontspec}
\setmainfont{DejaVu Serif}
\setsansfont{DejaVu Sans}
\setmonofont{DejaVu Sans Mono}

% ── Language ─────────────────────────────────────────────────────────────────
\usepackage{polyglossia}
\setmainlanguage{russian}
\setotherlanguage{english}

% ── Page layout ──────────────────────────────────────────────────────────────
\usepackage[
  top=25mm, bottom=25mm,
  left=30mm, right=20mm,
  headheight=14pt
]{geometry}

% ── Typography ───────────────────────────────────────────────────────────────
\usepackage{setspace}
\onehalfspacing

% ── Colours ──────────────────────────────────────────────────────────────────
\usepackage[dvipsnames,table]{xcolor}
\definecolor{primary}{HTML}{89B4FA}
\definecolor{success}{HTML}{A6E3A1}
\definecolor{danger}{HTML}{F38BA8}
\definecolor{warning}{HTML}{FAB387}
\definecolor{headerbg}{HTML}{1E1E2E}

% ── Headings ─────────────────────────────────────────────────────────────────
\usepackage{titlesec}
\titleformat{\section}
  {\large\bfseries\color{headerbg}}
  {\thesection.}{0.5em}{}[\titlerule]
\titleformat{\subsection}
  {\normalsize\bfseries}
  {\thesubsection.}{0.5em}{}
\titleformat{\subsubsection}
  {\normalsize\itshape}
  {\thesubsubsection.}{0.5em}{}

% ── Tables ───────────────────────────────────────────────────────────────────
\usepackage{booktabs}
\usepackage{tabularx}
\usepackage{array}
\renewcommand{\arraystretch}{1.4}

% ── Lists ────────────────────────────────────────────────────────────────────
\usepackage{enumitem}
\setlist[itemize]{noitemsep, topsep=4pt}
\setlist[enumerate]{noitemsep, topsep=4pt}

% ── Callout boxes ────────────────────────────────────────────────────────────
\usepackage{tcolorbox}
\tcbuselibrary{skins,breakable}

\newtcolorbox{warningbox}{
  colback=warning!15, colframe=warning!70,
  fonttitle=\bfseries, title=Внимание,
  breakable, before upper={\setlength\parindent{0pt}}
}
\newtcolorbox{notebox}{
  colback=primary!10, colframe=primary!60,
  fonttitle=\bfseries, title=Примечание,
  breakable, before upper={\setlength\parindent{0pt}}
}

% ── Hyperlinks ────────────────────────────────────────────────────────────────
\usepackage[
  unicode,
  colorlinks=true,
  linkcolor=headerbg,
  urlcolor=primary,
  pdftitle={Руководство пользователя — Центр Восстановления CRM},
  pdfauthor={Центр Восстановления},
  pdfsubject={Руководство оператора}
]{hyperref}

% ── Header / footer ──────────────────────────────────────────────────────────
\usepackage{fancyhdr}
\pagestyle{fancy}
\fancyhf{}
\fancyhead[L]{\small\itshape Центр Восстановления CRM}
\fancyhead[R]{\small\itshape Руководство пользователя}
\fancyfoot[C]{\small\thepage}
\renewcommand{\headrulewidth}{0.4pt}

% ── Button macros ─────────────────────────────────────────────────────────────
\newcommand{\btn}[2][primary]{%
  \textcolor{white}{\colorbox{#1}{\small\bfseries\,#2\,}}%
}
\newcommand{\btnred}[1]{\btn[danger]{#1}}
\newcommand{\btngreen}[1]{\btn[success!70!black]{#1}}
\newcommand{\btnyellow}[1]{\btn[warning!80!black]{#1}}

% =============================================================================
\begin{document}

% ── Title page ───────────────────────────────────────────────────────────────
\begin{titlepage}
  \centering
  \vspace*{4cm}
  {\Huge\bfseries Центр Восстановления CRM\par}
  \vspace{1cm}
  {\LARGE Руководство пользователя\par}
  \vspace{0.5cm}
  {\large Версия~1.0\par}
  \vspace{3cm}
  \rule{0.6\textwidth}{1pt}\\[0.4cm]
  {\large Руководство оператора системы\\[4pt]
  для работы с клиентами, тренировочными\\[2pt]
  программами и финансовым учётом\par}
  \vfill
  {\small\today\par}
\end{titlepage}

% ── Table of contents ────────────────────────────────────────────────────────
\tableofcontents
\newpage

% =============================================================================
\section{Общий интерфейс}
% =============================================================================

\subsection{Запуск программы}

После запуска исполняемого файла программа открывается в оконном режиме
(рекомендуемый размер~--- 1100\,$\times$\,700~пикселей и более).
По умолчанию отображается раздел \textbf{«Головна / Dashboard»}.

Все данные хранятся \textbf{локально} на компьютере в файле базы данных
\texttt{crm.db}, расположенном рядом с исполняемым файлом.
Резервное копирование выполняется вручную~--- просто скопируйте файл
\texttt{crm.db} в безопасное место.

\subsection{Боковое меню навигации}

Слева расположена вертикальная панель навигации шириной~200~px.
Активный раздел выделяется светлым прямоугольным фоном.
Для перехода в нужный раздел нажмите на его название.

\medskip
\begin{center}
\begin{tabularx}{0.85\textwidth}{>{\bfseries}lX}
  \toprule
  Пункт меню & Назначение \\
  \midrule
  Головна / Dashboard &
    Сводка активных программ клиентов, отметка посещений
    и выполнения упражнений \\
  Клієнти / Clients &
    Управление карточками клиентов и назначение программ \\
  Програми / Programs &
    Создание и редактирование тренировочных программ и упражнений \\
  Тренери / Trainers &
    Управление базой тренеров \\
  Фінанси / Finances &
    Учёт платежей с разбивкой по месяцам \\
  Розклад / Schedule &
    Еженедельный тайм-слотовый журнал записей \\
  \bottomrule
\end{tabularx}
\end{center}

\subsection{Переключение языка интерфейса}

В самом низу боковой панели находится переключатель языка~---
таблетка с двумя кнопками:

\begin{itemize}
  \item \textbf{EN}~--- английский интерфейс
  \item \textbf{UA}~--- украинский интерфейс (значение по умолчанию)
\end{itemize}

Нажмите нужную кнопку. Переключение происходит мгновенно и затрагивает
только текст интерфейса; данные в базе не изменяются.

% =============================================================================
\section{Раздел «Головна / Dashboard»}
% =============================================================================

\subsection{Назначение}

Раздел «Главная» служит рабочим столом тренера на каждый день.
Он отображает все активные программы клиентов: имя клиента, название
программы, тренера, дату начала. Здесь же отмечаются пройденные недели
и выполненные упражнения, а также формируется и печатается PDF-отчёт.

\subsection{Кнопки печати и сохранения}

В правом верхнем углу расположены две кнопки:

\begin{warningbox}
  Кнопки становятся активными \emph{только после того, как вы нажмёте
  на карточку программы} в списке. Пока ни одна карточка не выбрана,
  кнопки затемнены и не реагируют на нажатие.
\end{warningbox}

\bigskip
\begin{center}
\begin{tabularx}{0.9\textwidth}{>{\bfseries}lX}
  \toprule
  Кнопка & Действие \\
  \midrule
  \btn{Зберегти PDF / Save PDF} &
    Открывает системное диалоговое окно сохранения файла.
    Выберите папку и имя, нажмите «Сохранить».
    Программа создаст PDF-отчёт \emph{только по выбранной программе}. \\[4pt]
  \btn[success!70!black]{Друк / Print} &
    Открывает системный диалог печати.
    Отчёт содержит данные \emph{только выбранной программы}:
    список недель (check-ins) и таблицу упражнений. \\
  \bottomrule
\end{tabularx}
\end{center}

\subsection{Карточки активных программ}

Каждая карточка отображает:
\begin{itemize}
  \item Имя клиента и название программы
  \item Имя тренера
  \item Дату начала программы
\end{itemize}

\textbf{Нажмите на карточку}, чтобы её развернуть.
Повторное нажатие сворачивает карточку.
При раскрытии ниже появляются два блока: «Відвідування» и «Вправи».

\subsection{Блок «Відвідування / Check-ins»}

Список недель программы. Каждая строка содержит:
\begin{itemize}
  \item Номер недели (Тиж./Wk)
  \item Дату начала недели
  \item Чекбокс выполнения (\textbf{[+]} выполнено~/ \textbf{[ ]} не выполнено)
\end{itemize}

Нажмите на чекбокс, чтобы отметить неделю как завершённую.
Повторное нажатие снимает отметку. Если недели не запланированы~---
отображается «Відвідування не заплановані / No check-ins scheduled».

\subsection{Блок «Вправи / Exercises»}

Список упражнений текущей программы. Каждая строка показывает:
\begin{itemize}
  \item Название упражнения
  \item Схему нагрузки: подходы\,$\times$\,повторения,
    например \texttt{3x10}
  \item Рабочий вес, если задан: \texttt{@60kg}
  \item Заметки к упражнению
  \item Чекбокс выполнения
\end{itemize}

Если упражнений нет~--- отображается «У цій програмі немає вправ /
No exercises in this program».

% =============================================================================
\section{Раздел «Клієнти / Clients»}
% =============================================================================

\subsection{Назначение}

Хранение базы клиентов с контактными данными, привязкой к тренеру
и историей назначенных программ.

\subsection{Кнопка «Додати клієнта / Add Client»}

Расположена в правом верхнем углу. При нажатии открывается форма
добавления нового клиента; кнопка меняет цвет на красный и текст
на \btnred{Скасувати / Cancel}. Нажмите повторно, чтобы закрыть
форму \emph{без сохранения}.

\subsection{Форма добавления / редактирования клиента}

\begin{center}
\begin{tabularx}{0.9\textwidth}{>{\bfseries}lX}
  \toprule
  Поле & Описание \\
  \midrule
  Ім'я / Name &
    Полное имя клиента. Текстовое поле, обязательно для заполнения. \\
  Телефон / Phone &
    Номер телефона в произвольном формате. \\
  Ел. пошта / Email &
    Электронная почта. \\
  Нотатки / Notes &
    Любые заметки (диагноз, особенности здоровья, пожелания). \\
  Тренер / Trainer &
    Пикер выбора тренера: кнопки \textbf{<} и \textbf{>} листают
    список тренеров. Значение «Немає / None»~--- тренер не назначен. \\
  \bottomrule
\end{tabularx}
\end{center}

\medskip
\noindent После заполнения нажмите \btn{Зберегти / Save}.
При редактировании существующего клиента кнопка меняет текст
на \btn{Оновити / Update}.

\subsection{Список клиентов}

Таблица с колонками: \textit{Ім'я, Телефон, Ел.~пошта, Тренер, Дії}.

В колонке \textbf{«Дії»} каждой строки:
\begin{itemize}
  \item \btn{Змінити / Edit}~--- открывает форму с предзаполненными
    данными выбранного клиента для редактирования.
  \item \btnred{Видалити / Delete}~--- немедленно и \textbf{без
    подтверждения} удаляет клиента из базы.
    Будьте внимательны.
\end{itemize}

\begin{warningbox}
  Удаление клиента не удаляет его финансовые записи и записи
  в расписании. Они останутся в базе, но без привязки к имени.
\end{warningbox}

Нажмите \textbf{на строку} (не на кнопки), чтобы выбрать клиента~---
справа откроется панель его программ.

\subsection{Правая панель~--- программы клиента}

Панель появляется при выбранном клиенте.

\subsubsection{Назначение программы}

\begin{enumerate}
  \item С помощью кнопок \textbf{<} / \textbf{>} выберите программу
    из списка.
  \item Нажмите \btngreen{Призначити / Assign}.
    Программа стартует с текущей даты и получает статус
    \textbf{«Активна»}.
\end{enumerate}

\subsubsection{Список назначенных программ}

Каждая строка содержит цветной индикатор статуса, название программы,
дату начала и кнопки управления.

\medskip
\begin{center}
\begin{tabular}{llp{6.2cm}}
  \toprule
  Цвет & Статус & Доступные кнопки \\
  \midrule
  \textcolor{success!70!black}{\rule{10pt}{10pt}}~Зелёный &
    Активна &
    \btnyellow{Пауза / Pause}~\quad~\btn{Завершити / Complete} \\[2pt]
  \textcolor{warning}{\rule{10pt}{10pt}}~Жёлтый &
    На паузі &
    \btngreen{Продовжити / Resume}~\quad~\btnred{Скасувати / Cancel} \\[2pt]
  \textcolor{primary!80!black}{\rule{10pt}{10pt}}~Синий &
    Завершена &
    (нет, статус финальный) \\[2pt]
  \textcolor{gray}{\rule{10pt}{10pt}}~Серый &
    Скасована &
    (нет, статус финальный) \\
  \bottomrule
\end{tabular}
\end{center}

% =============================================================================
\section{Раздел «Програми / Programs»}
% =============================================================================

\subsection{Назначение}

Создание шаблонов тренировочных программ с набором упражнений.
Программы затем назначаются клиентам из раздела «Клієнти».

\begin{notebox}
  Изменение программы отражается у всех клиентов, которым она назначена.
  Планируйте шаблоны заранее.
\end{notebox}

\subsection{Форма добавления / редактирования программы}

\begin{center}
\begin{tabularx}{0.9\textwidth}{>{\bfseries}lX}
  \toprule
  Поле & Описание \\
  \midrule
  Назва / Name &
    Название программы (произвольный текст). \\
  Тривалість (тижні) / Duration (weeks) &
    Счётчик от \textbf{1} до \textbf{52}. Нажимайте
    \textbf{+} / \textbf{-} для изменения. Определяет количество
    недель (check-in строк) в разделе «Головна». \\
  Опис / Description &
    Описание программы, видимое в таблице списка. \\
  \bottomrule
\end{tabularx}
\end{center}

\subsection{Список программ}

Таблица: \textit{Назва, Опис, Тижні, Дії}.
Нажмите \textbf{на строку программы}, чтобы открыть справа панель
упражнений.

\subsection{Правая панель~--- упражнения программы}

\subsubsection{Форма упражнения}

\begin{center}
\begin{tabularx}{0.9\textwidth}{>{\bfseries}lX}
  \toprule
  Поле & Описание \\
  \midrule
  Назва / Name &
    Название упражнения (например: «Жим лёжа», «Присед»). \\
  Підходи / Sets &
    Счётчик от \textbf{1} до \textbf{20}~---
    количество подходов. \\
  Повтори / Reps &
    Счётчик от \textbf{1} до \textbf{100}~---
    количество повторений. \\
  Вага (кг) / Weight (kg) &
    Счётчик от \textbf{0} до \textbf{999}~---
    рабочий вес в кг. Значение \textbf{0}~---
    упражнение без отягощения (вес тела). \\
  Нотатки / Notes &
    Необязательные заметки (темп, техника, ограничения). \\
  \bottomrule
\end{tabularx}
\end{center}

\subsubsection{Список упражнений}

Каждая строка отображает название, нагрузку (\texttt{3x10 @60kg})
и заметки. Кнопки: \btn{Змінити / Edit} и \btnred{Видалити / Delete}.

% =============================================================================
\section{Раздел «Тренери / Trainers»}
% =============================================================================

\subsection{Назначение}

Ведение базы тренеров. Тренеры привязываются к клиентам.

\subsection{Форма тренера}

\begin{center}
\begin{tabularx}{0.9\textwidth}{>{\bfseries}lX}
  \toprule
  Поле & Описание \\
  \midrule
  Ім'я / Name &
    Полное имя тренера. \\
  Телефон / Phone &
    Рабочий номер телефона. \\
  Спеціалізація / Specialization &
    Специализация, например «Фізіотерапія»,
    «Лікувальна фізкультура», «Силові тренування». \\
  \bottomrule
\end{tabularx}
\end{center}

\subsection{Список тренеров}

Таблица: \textit{Ім'я, Телефон, Спеціалізація, Дії}.
Кнопки: \btn{Змінити / Edit} и \btnred{Видалити / Delete}.

\begin{warningbox}
  Удаление тренера \emph{не} удаляет клиентов, привязанных к нему.
  У этих клиентов поле «Тренер» станет пустым («Немає»).
\end{warningbox}

% =============================================================================
\section{Раздел «Фінанси / Finances»}
% =============================================================================

\subsection{Назначение}

Учёт всех платежей клиентов с разбивкой по месяцам.
Автоматически показывает итоговую сумму за выбранный месяц.

\subsection{Навигация по месяцам}

В заголовке раздела:
\begin{itemize}
  \item \textbf{<}~--- переход к предыдущему месяцу
  \item Дисплей в формате \texttt{ГГГГ-ММ} (например: \texttt{2026-02})
  \item \textbf{>}~--- переход к следующему месяцу
\end{itemize}

Список и итоговая сумма обновляются автоматически при смене месяца.

\subsection{Итоговая сумма}

Зелёная плашка справа от навигации:
\textbf{Всього / Total: СУММА грн/UAH}~---
суммарный доход за текущий выбранный месяц.

\subsection{Форма добавления / редактирования визита}

Кнопка \btn{Додати візит / Add Visit} открывает форму.
При повторном нажатии форма закрывается без сохранения.

\bigskip
\begin{center}
\begin{tabularx}{0.9\textwidth}{>{\bfseries}lX}
  \toprule
  Поле & Описание \\
  \midrule
  Клієнт / Client &
    Пикер \textbf{<} / \textbf{>}~--- выбор клиента из базы. \\
  Дата / Date &
    Дата в формате \texttt{ГГГГ-ММ-ДД} (например: \texttt{2026-02-25}).
    При открытии формы устанавливается первое число текущего месяца. \\
  Сума / Amount &
    Сумма платежа в гривнях. Вводите цифры (допустимо с десятичными:
    \texttt{500.00}). \\
  Нотатки / Notes &
    Комментарий (вид услуги, источник платежа и т.\,д.). \\
  \bottomrule
\end{tabularx}
\end{center}

\begin{notebox}
  Записи из раздела «Розклад» попадают сюда автоматически~---
  при нажатии кнопки «Прийшов / Attended» с указанной суммой.
  Такие записи будут иметь нотатку «Запис (зустріч) / Appointment».
\end{notebox}

\subsection{Список визитов}

Таблица: \textit{Дата, Клієнт, Сума, Нотатки, Дії}.
Суммы отображаются зелёным цветом с единицей валюты (грн / UAH).
Кнопки: \btn{Змінити / Edit} и \btnred{Видалити / Delete}.

% =============================================================================
\section{Раздел «Розклад / Schedule»}
% =============================================================================

\subsection{Назначение}

Тижневий розклад~--- еженедельная сетка от \textbf{06:00}
до \textbf{21:00}, разделённая на 7 дневных колонок.
Позволяет бронировать время для клиентов, просматривать загрузку
на неделю и отмечать фактическую посещаемость.

\subsection{Навигация по неделям}

\begin{itemize}
  \item \textbf{<}~--- предыдущая неделя
  \item Дисплей текущей недели (например: \texttt{Feb 23 – Mar 1, 2026})
  \item \textbf{>}~--- следующая неделя
\end{itemize}

\subsection{Форма добавления / редактирования записи}

Кнопка \btn{Додати запис / Add Appointment} открывает форму.
При повторном нажатии форма закрывается без сохранения.

\bigskip
\begin{center}
\begin{tabularx}{0.9\textwidth}{>{\bfseries}lX}
  \toprule
  Поле & Описание \\
  \midrule
  Клієнт / Client &
    Пикер \textbf{<} / \textbf{>}~--- листать список клиентов. \\
  Дата / Date &
    Дата в формате \texttt{ГГГГ-ММ-ДД}. При клике по ячейке сетки
    заполняется автоматически. \\
  Час / Time &
    Два пикера: \textbf{часы} (диапазон 06--21, шаг 1~час) и
    \textbf{минуты} (шаг 5~мин: 00, 05, 10, \ldots, 55).
    Кнопки \textbf{<} / \textbf{>} изменяют значение. \\
  Тривалість (хв) / Duration (min) &
    Длительность записи: кнопки \textbf{<} / \textbf{>} с шагом
    15~минут (диапазон 15--240~мин). \\
  Нотатки / Notes &
    Произвольные заметки. \\
  \bottomrule
\end{tabularx}
\end{center}

Кнопки формы (нижняя строка):
\begin{itemize}
  \item \btn{Зберегти / Save}~--- сохраняет новую запись и закрывает форму.
  \item \btn{Оновити / Update}~--- сохраняет изменения существующей записи
    (заменяет «Зберегти» при редактировании).
  \item \btnred{Видалити / Delete}~--- удаляет запись
    (видна только при редактировании существующей).
\end{itemize}

\subsection{Блок отметки посещаемости}

\begin{notebox}
  Этот блок отображается \emph{только} при редактировании существующей
  записи (когда вы нажали на карточку в сетке расписания).
\end{notebox}

\bigskip
\begin{center}
\begin{tabularx}{0.9\textwidth}{>{\bfseries}lX}
  \toprule
  Элемент & Действие \\
  \midrule
  Оплата / Payment &
    Текстовое поле суммы. Введите сумму оплаты за визит.
    Если оставить \texttt{0}~--- финансовая запись не создаётся. \\[4pt]
  \btngreen{Прийшов / Attended} &
    Помечает запись как «клиент пришёл».
    Карточка в сетке становится \textbf{зелёной}.
    Если сумма~$>$~0~--- автоматически создаётся запись в
    разделе «Фінанси» на указанную сумму. \\[4pt]
  \btnred{Не прийшов / No-show} &
    Помечает запись как «клиент не пришёл».
    Карточка в сетке становится \textbf{красной}.
    Финансовая запись не создаётся. \\
  \bottomrule
\end{tabularx}
\end{center}

\subsection{Временная сетка}

При открытии раздела сетка автоматически прокручивается к~\textbf{08:00}.
Прокрутка выполняется колесом мыши или трекпадом.

\subsubsection{Шкала времени (левая колонка)}

Отметки слева от сетки:
\begin{itemize}
  \item \textbf{Метка часа}~--- подпись вида \texttt{08:00}, \texttt{09:00}~\ldots
  \item \textbf{Засечки 5 минут}~--- короткие линии
  \item \textbf{Засечки 15 минут}~--- средние линии
  \item \textbf{Засечка 30 минут}~--- длинная линия
\end{itemize}

\subsubsection{Заголовки дней}

Текущий день выделяется \textbf{синим прямоугольником}.
Остальные дни~--- серым фоном.

\subsubsection{Создание записи кликом по ячейке сетки}

\begin{enumerate}
  \item Прокрутите сетку до нужного времени.
  \item Нажмите на ячейку нужного дня и часа.
    Ячейка подсвечивается, форма открывается с предзаполненными
    датой и временем.
  \item Выберите клиента, скорректируйте время и длительность при
    необходимости, добавьте заметки.
  \item Нажмите \btn{Зберегти / Save}. Карточка появится в сетке.
\end{enumerate}

\subsubsection{Карточки записей в сетке}

\begin{center}
\begin{tabular}{llp{7cm}}
  \toprule
  Цвет карточки & Статус & Значение \\
  \midrule
  \textcolor{primary}{\rule{12pt}{12pt}}~Синий &
    pending & Запись создана, посещаемость не отмечена \\[2pt]
  \textcolor{success!70!black}{\rule{12pt}{12pt}}~Зелёный &
    attended & Клиент пришёл \\[2pt]
  \textcolor{danger}{\rule{12pt}{12pt}}~Красный &
    no-show & Клиент не пришёл \\
  \bottomrule
\end{tabular}
\end{center}

На карточке отображается имя клиента; при длительности 45~минут
и более~--- также время начала записи.
Нажмите на карточку, чтобы открыть форму редактирования.

% =============================================================================
\section{Порядок работы и типовые сценарии}
% =============================================================================

\subsection{Первоначальная настройка системы}

\begin{enumerate}
  \item \textbf{Тренери}~--- добавьте всех сотрудников-тренеров.
  \item \textbf{Програми}~--- создайте шаблоны программ. Для каждой
    программы добавьте упражнения с нагрузкой (подходы, повторения, вес).
  \item \textbf{Клієнти}~--- добавьте клиентов, привяжите тренеров.
  \item \textbf{Клієнти, правая панель}~--- назначьте клиентам программы.
  \item \textbf{Головна}~--- убедитесь, что карточки клиентов появились
    в разделе «Главная».
\end{enumerate}

\subsection{Ежедневная работа тренера}

\begin{enumerate}
  \item Откройте \textbf{Розклад}~--- просмотрите записи на сегодня.
  \item После тренировки нажмите на карточку клиента в расписании.
  \item Введите сумму оплаты в поле «Оплата / Payment».
  \item Нажмите \btngreen{Прийшов / Attended}~---
    посещение зафиксировано, платёж автоматически внесён в
    \textbf{Фінанси}.
  \item Если клиент не пришёл~--- нажмите
    \btnred{Не прийшов / No-show}.
  \item В разделе \textbf{Головна} откройте карточку нужной программы
    и отметьте выполненные упражнения.
\end{enumerate}

\subsection{Составление отчёта по программе}

\begin{enumerate}
  \item Перейдите в \textbf{Головна}.
  \item Нажмите на карточку нужного клиента~--- она выделится.
  \item Нажмите \btn{Зберегти PDF / Save PDF} или
    \btn[success!70!black]{Друк / Print}.
  \item Выберите путь сохранения (для PDF) или принтер (для печати).
\end{enumerate}

\subsection{Общие правила и ограничения}

\begin{itemize}
  \item Все данные сохраняются \textbf{автоматически} при нажатии
    «Зберегти». Отдельного сохранения нет.
  \item Кнопки \btnred{Видалити / Delete} \textbf{не запрашивают
    подтверждения}. Удалённые данные восстановить невозможно.
  \item Дата вводится вручную в формате \texttt{ГГГГ-ММ-ДД}.
  \item Для резервной копии периодически копируйте файл \texttt{crm.db}.
\end{itemize}

\vfill
\begin{center}
  \rule{0.5\textwidth}{0.4pt}\\[4pt]
  {\small Версия~1.0 \quad \today}
\end{center}

\end{document}
